% (User input window)

\subsection{纤维谱的行列式势：多重度算子与反射投影的同一温度读出}\label{subsec:op-algebra-fiber-determinant-potentials}
命题 \ref{prop:fold-fiber-multiplicity-trace-moments} 已表明：折叠纤维多重度 \(\{d_m(x)\}_{x\in X_m}\) 可被提升为稳定层上的对角谱算子 \(\mathcal{F}_m\)，其迹矩 \(\Tr(\mathcal{F}_m^q)\) 精确等于碰撞矩 \(S_q(m)\)。
本小节记录与迹矩互补的另一条谱函子：由同一谱数据 \(\{d_m(x)\}\) 生成的行列式势（determinant potential），它在制度上把“微观体积、边界体积与纤维谱”压缩为同一配分函数在不同温度下的读出。

\begin{proposition}[纤维谱配分函数：\texorpdfstring{$Z_m(t)=\det(I+t\mathcal{F}_m)$}{Zm(t)=det(I+tF)} 的闭式与两端读出]\label{prop:op-algebra-fiber-spectrum-determinant-partition}
\leanverified{paper\_op\_algebra\_fiber\_spectrum\_determinant\_partition}
在定义 \ref{def:fold-fiber-multiplicity-operator} 的记号下，对任意 \(t\ge 0\) 定义行列式势
\[
Z_m(t):=\det(I+t\mathcal{F}_m).
\]
则有完全闭式
\[
\boxed{\;
Z_m(t)=\prod_{x\in X_m}\bigl(1+t\,d_m(x)\bigr)\;}
\]
并且其小温度/高温极限满足
\[
\boxed{\;
\lim_{t\downarrow 0}\frac{Z_m(t)-1}{t}=\Tr(\mathcal{F}_m)=\sum_{x\in X_m}d_m(x)=|\Omega_m|=2^m\;}
\]
以及
\[
\boxed{\;
\lim_{t\to\infty}\frac{\dd}{\dd(\log t)}\log Z_m(t)
=\lim_{t\to\infty}\sum_{x\in X_m}\frac{t\,d_m(x)}{1+t\,d_m(x)}
=|X_m|\;}
\]
从而
\[
\boxed{\;
\lim_{t\to\infty}\Bigl(\log Z_m(t)-|X_m|\log t\Bigr)=\sum_{x\in X_m}\log d_m(x)\;}
\]
给出一条把边界体积 \(|X_m|\) 与纤维谱 \(\{d_m(x)\}\) 同时封装进单一行列式势的严格插值读出。
\end{proposition}
\begin{proof}
由命题 \ref{prop:fold-fiber-multiplicity-trace-moments}，\(\mathcal{F}_m\) 在标准基下为对角算子且本征值为 \(\{d_m(x)\}_{x\in X_m}\)，故
\(
Z_m(t)=\prod_x(1+t d_m(x))
\)。
对 \(t\downarrow 0\) 作一阶展开即得
\(
Z_m(t)=1+t\Tr(\mathcal{F}_m)+O(t^2)
\)；
而 \(\sum_x d_m(x)=|\Omega_m|\) 为纤维基数求和恒等式。
\par
对导数，直接由
\(
\log Z_m(t)=\sum_x \log(1+t d_m(x))
\)
可得
\(
\frac{\dd}{\dd(\log t)}\log Z_m(t)=\sum_x \frac{t d_m(x)}{1+t d_m(x)}
\to |X_m|
\)（逐项极限）。
最后一条由 \(\log(1+t d)=\log t+\log(d+1/t)\) 并令 \(t\to\infty\) 得到。证毕。
\end{proof}

\begin{corollary}[纤维反射子的行列式势：\texorpdfstring{$\det(I+tK_m)=(1+t)^{|X_m|}$}{det(I+tK)=(1+t)^|X|}]\label{cor:op-algebra-fiber-reflector-determinant-potential}
把纤维反射子 \(K_m\) 视为 \(\CC^{\Omega_m}\) 上的线性算子（定理 \ref{thm:pom-fiber-reflector} 与推论 \ref{cor:pom-Km-spectrum}），则对任意 \(t>-1\) 有
\[
\boxed{\;
\det(I+tK_m)=(1+t)^{|X_m|},\qquad
\frac{1}{m}\log\det(I+tK_m)=\frac{|X_m|}{m}\log(1+t)\;}
\]
从而边界体积 \(|X_m|\) 的指数增长率可由任意固定 \(t\neq 0\) 的归一化行列式势无损读出。
\end{corollary}
\begin{proof}
由推论 \ref{cor:pom-Km-spectrum}，\(K_m\) 的谱仅由 \(\{0,1\}\) 构成且特征值 \(1\) 的重数为 \(|X_m|\)。
因此 \(I+tK_m\) 的特征值为 \(1+t\)（重数 \(|X_m|\)）与 \(1\)（重数 \(2^m-|X_m|\)），取行列式得结论。证毕。
\end{proof}

\begin{proposition}[行列式势的符号翻转与容量斜率奇偶性]\label{prop:op-algebra-fiber-determinant-sign-capacity-parity}
令 \(Z_m(t)=\prod_{x\in X_m}(1+t\,d_m(x))\) 如命题 \ref{prop:op-algebra-fiber-spectrum-determinant-partition}。
对 \(T>0\) 且 \(T\neq d_m(x)\)（对所有 \(x\)）定义
\[
S_m(T):=\mathrm{sgn}\,Z_m(-1/T)\in\{\pm 1\}.
\]
则在每个正则点 \(T\) 有严格恒等式
\[
\boxed{\;
S_m(T)=(-1)^{\#\{x\in X_m:\ d_m(x)>T\}}
=(-1)^{\mathcal{C}_m'(T)},\;}
\]
其中 \(\mathcal{C}_m(T)=\sum_x\min(d_m(x),T)\) 及其导数见命题 \ref{prop:fold-oracle-capacity-slope-jumps-determine-histogram}。
因此 \(S_m(T)\) 仅在 \(T\) 穿越某个出现过的纤维大小 \(d\) 时发生符号翻转，且翻转当且仅当该大小的出现次数 \(N_{m,d}\) 为奇数。
\end{proposition}
\begin{proof}
由乘积分解，
\[
Z_m(-1/T)=\prod_{x\in X_m}\Bigl(1-\frac{d_m(x)}{T}\Bigr).
\]
在正则点 \(T\) 上，每一因子符号为 \(\mathrm{sgn}(T-d_m(x))\)，故
\[
\mathrm{sgn}\,Z_m(-1/T)=(-1)^{\#\{x:\ d_m(x)>T\}}.
\]
再由命题 \ref{prop:fold-oracle-capacity-slope-jumps-determine-histogram} 得 \(\mathcal{C}_m'(T)=\#\{x:\ d_m(x)>T\}\)，合并得到第一式。
当 \(T\) 穿越整数 \(d\) 时，该计数恰减少 \(N_{m,d}\)，故符号翻转当且仅当 \(N_{m,d}\) 为奇数。证毕。
\end{proof}
